\lecture{3}{3 March.}{}

\begin{definition}
    Given events \(E_1, E_2, \dots \subseteq S\), 
    \begin{itemize}
        \item \(\bigcup_{i=1}^{\infty} E_i = \left\{ x \in S : \text{there exists } i \text{ with } x \in E_i   \right\}  \). 
        \item \(\bigcap_{i=1}^{\infty} E_i = \left\{ x \in S: x \in E_i \text{ for all } i  \right\}  \). 
        \item \(E_1, E_2, \dots \) are mutually exclusive (or pairwise disjoint) if \(E_i \cap E_j = \varnothing \) for all \(i \neq j\).     
        \item Given \(E \subseteq S\), the complement \(E^C = \overline{E} = \left\{ x \in S: x \notin E \right\} \). 
        \item Given \(E, F \subseteq S\), \(F \setminus E = \left\{ x \in F: x \notin E \right\} = F \cap E^C \). 
        \item \(E \Delta F = (E \setminus F) \cup (F \setminus E)\).     
    \end{itemize} 
\end{definition}

\begin{proposition}[De Morgan's law]
    Given sets \(E_1, E_2, \dots \subseteq S\), 
    \[
        \left( \bigcup_{i=1}^{\infty} E_i  \right)^C = \bigcap_{i=1}^{\infty} E_i^C \text{ and } \left( \bigcap_{i=1}^{\infty} E_i  \right)^C = \bigcup_{i=1}^{\infty} E_i^C.     
    \] 
\end{proposition}

\begin{proof}
    \begin{align*}
        \left( \bigcup_{i=1}^{\infty} E_i  \right)^C &= \left\{ x \in S : x \notin \bigcup_{i=1}^{\infty} E_i  \right\} = \left\{ x \in S: \not\exists i \text{ s.t. } x \in E_i  \right\} \\
        &= \left\{ x \in S: \forall i, x \notin E_i \right\} = \left\{ x \in S: \forall i, x \in E_i^C \right\} = \bigcap_{i=1}^{\infty} E_i^C .  \\
        \left( \bigcap_{i=1}^{\infty} E_i  \right)^C &= \left( \bigcap_{i=1}^{\infty} \left( E_i^C \right)^C   \right)^C = \left( \left( \bigcup_{i=1}^{\infty} E_i^C  \right)^C \right)^C = \left( \bigcup_{i=1}^{\infty} E_i^C  \right).         
    \end{align*}
\end{proof}

\begin{definition}
    Given a sample space \(S\), a probability measure is a function \(\mathbb{P} : 2^S \to \mathbb{R} \) satisfying 
    \begin{itemize}
        \item [(i)] for every event \(E\), we have \(0 \le \mathbb{P} (E) \le 1\). 
        \item [(ii)] \(\mathbb{P} (S) = 1\). 
        \item [(iii)] (countable additivity) for every sequence of mutually exclusive events \(E_1, E_2, \dots \)
        \[
            \mathbb{P} \left( \bigcup_{i=1}^{\infty} E_i  \right) = \sum_{i=1}^{\infty} \mathbb{P} (E_i).  
        \]
    \end{itemize}   
\end{definition}

\begin{definition}
    Given a sample space \(S\), and a probability measure \(\mathbb{P} :2^S \to \mathbb{R} \), we call \((S, \mathbb{P} )\) a probability space.
\end{definition}

\begin{proposition}
    \(\mathbb{P} (\varnothing ) = 0\). 
\end{proposition}  

\begin{proof}
    Consider the following sequence of events \(E_1 = S\), \(E_i = \varnothing \) for all \(i \ge 2\). Observe that those are mutually exclusive: 
\[
    \forall i \neq j, E_i \cap E_j = \varnothing
\]   
since one of \(E_i, E_j\) must be empty. Thus, 
\[
    \mathbb{P} \left( \bigcup_{i=1}^{\infty} E_i  \right) = \sum_{i=1}^{\infty} \mathbb{P} (E_i),  
\] 
and since 
\[
    \mathbb{P} \left( \bigcup_{i=1}^{\infty} E_i  \right)  = \mathbb{P} (S) = 1,
\]
we know 
\[
    1 = \sum_{i=1}^{\infty} \mathbb{P} (E_i) = \mathbb{P} (S) + \sum_{i=2}^{\infty} \mathbb{P} (\varnothing ) = 1 + \sum_{i=2}^{\infty} \mathbb{P} (\varnothing ).   
\]
This gives \(\mathbb{P} (\varnothing ) = 0\).
\end{proof}

\begin{proposition}
    If \(E_1, E_2, \dots , E_n\) is a finite sequence of mutually exclusive events, then 
    \[
        \mathbb{P} \left( \bigcup_{i=1}^{n} E_i  \right) = \sum_{i=1}^n \mathbb{P} (E_i).  
    \] 
\end{proposition} 

\begin{proof}
    Consider the sequence 
    \[
        E_1, E_2, \dots , E_n, \varnothing , \varnothing , \dots,
    \]
    then we have 
    \[
         \mathbb{P} \left( \bigcup_{i=1}^{n} E_i  \right) = \mathbb{P} \left( \bigcup_{i=1}^{\infty} E_i  \right) = \sum_{i=1}^{\infty } \mathbb{P} (E_i) = \sum_{i=1}^n \mathbb{P} (E_i).   
    \]
\end{proof}

\begin{proposition}
    For any event \(E\), 
    \[
        \mathbb{P} \left( E^C \right) = 1 - \mathbb{P} (E). 
    \] 
\end{proposition}

\begin{proof}
    Consider the finite sequence 
    \[
        E_1 = E \text{ and } E_2 = E^C, 
    \]
    then 
    \[
         \sum_{i=1}^2 \mathbb{P} (E_i) = \mathbb{P} \left( \bigcup_{i=1}^{2} E_i  \right) = \mathbb{P} (S) = 1, 
    \]
    which gives 
    \[
        \mathbb{P} (E_1) = 1 - \mathbb{P} \left( E_1^C \right). 
    \]
\end{proof}

\begin{proposition}[Monotonity]
    If \(E \subseteq F\) are events, then 
    \[
        \mathbb{P} (E) \le \mathbb{P} (F).
    \] 
\end{proposition}

\begin{proof}
    Note that \(F = E \cup \left( F \setminus E \right) \), so 
    \[
        \mathbb{P} (F) = \mathbb{P} (E) + \mathbb{P} (F \setminus E) \ge \mathbb{P} (E).
    \] 
\end{proof}

\begin{eg}
    We have an unfair six-sided die, which rolls \(1\) with probability \(0.1\), a \(2\) with probability \(0.2\), a \(3\) with probability \(0.25\), a \(4\) with probability \(0.15\), and a \(5\) with probability \(0.17\). Then, 
    \begin{itemize}
        \item [(a)] What is the probability of rolling a \(6\)? 
        \item [(b)] What is the probability of rolling a prime? 
    \end{itemize}         
\end{eg}
\begin{explanation}
    To define the probability over a countable sample space \(S = \left\{ 1, 2, 3, 4, 5, 6 \right\} \), it is enough to define the probability of individual outcomes \(\left\{ x \right\} \subseteq S \). Since for any \(E \subseteq S\), we know \(E = \bigcup_{x \in E}\left\{ x \right\}  \) and thus \(\mathbb{P} (E) = \sum_{x \in E} \mathbb{P} \left( \left\{ x \right\}  \right)  \). Hence, we know 
    \[
        \mathbb{P} (\left\{ 6 \right\} ) = \mathbb{P} \left( \left\{ 1,2,3,4,5 \right\}^C  \right) = 1 - \sum_{i=1}^5 \mathbb{P} \left( \left\{ i \right\}  \right) = 0.13.    
    \]
    Also, 
    \[
        \mathbb{P} (\text{rolling a prime} ) = 0.62.
    \]   
\end{explanation}

\begin{eg}
    Any uniform probability measure satisfies the axiom.
\end{eg}
\begin{explanation}
    Let \(S\) be a finite non-empty sample space. We define the uniform probability measure as 
    \[
        \mathbb{P} (E) = \frac{\vert E \vert }{\vert S \vert }.
    \]  
    Fix any \(E \subseteq S\), \(0 \le \vert E \vert \le \vert S \vert  \), so 
    \[
        0 \le \frac{\vert E \vert }{\vert S \vert } \le 1.
    \]
    Also, \(\mathbb{P} (S) = \frac{\vert S \vert }{\vert S \vert } = 1\). If \(E_1, E_2, \dots \) are mutually exclusive, then 
    \[
        \mathbb{P} \left( \bigcup_{i=1}^{\infty} E_i  \right) = \frac{\left\vert \bigcup_{i=1}^{\infty} E_i  \right\vert }{\vert S \vert } = \frac{\sum_{i=1}^{\infty} \vert E_i \vert  }{\vert S \vert } = \sum_{i=1}^{\infty} \frac{\vert E_i \vert }{\vert S \vert } = \sum_{i=1}^{\infty} \mathbb{P} (E_i).  
    \]  
    Thus, any uniform probability measure satisfies the axiom.
\end{explanation}

\section{Inclusion-Exclusion Principle}

\begin{proposition}
    For any two events \(E, F \subseteq S\), 
    \[
        \mathbb{P} (E \cup F) = \mathbb{P} (E) + \mathbb{P} (F) - \mathbb{P} (E \cap F).
    \] 
\end{proposition}
\begin{proof}
    Consider the mutually exclusive events \(E \cap F, E \setminus F, F \setminus E\), we know 
    \[
        \left( E \cap F \right) \cup \left( E \setminus F \right) \cup \left( F \setminus E \right) = E \cup F.   
    \] 
    Also, 
    \[
        (E \cap F) \cup (E \setminus F) = E \text{ and } (E \cap F) \cup (F \setminus E) = F. 
    \]
    Hence, 
    \begin{align*}
        \mathbb{P} (E \cup F) &= \mathbb{P} (E \cap F) + \mathbb{P} (E \setminus F) + \mathbb{P} (F \setminus E) \\
        &= \mathbb{P} (E \cap F) + \mathbb{P} (E \setminus F) + \mathbb{P} (F \setminus E) + \mathbb{P} (E \cap F) - \mathbb{P} (E \cap F) \\
        &= \mathbb{P} (E) + \mathbb{P} (F) - \mathbb{P} (E \cap F).
    \end{align*}
\end{proof}

\begin{corollary}
    For any \(E, F \subseteq S\), we know 
    \[
        \mathbb{P} (E \cup F) \le \mathbb{P} (E) + \mathbb{P} (F).
    \] 
\end{corollary}